\section{The Structural Cognition paradigm and Self-Aware Steel Structures}\label{sec:paradigm}

Conventional structural health monitoring (SHM) is defined as the process of implementing a damage detection strategy for engineering infrastructure: observing a structure over time through periodically spaced response measurements, extracting damage-sensitive features, and statistically analysing those features to determine the current state of system health \citep{farrar2013shm}. In its mature, deployed form this process is overwhelmingly \textit{extrospective}: instrumentation records strains, accelerations, temperatures and displacements, transmits them to a data acquisition system, and stores them for later interrogation by an engineer or an offline algorithm. The structure itself remains a passive load-resisting object, and the cognitive act of converting measurements into an assessment of integrity occurs elsewhere, asynchronously, and usually under human supervision. This is consistent with Axiom IVa of SHM, that sensors cannot measure damage, which establishes that feature extraction and statistical classification are always required to convert sensor data into damage information \citep{farrar2013shm,farrar2007introduction}. Conventional practice satisfies this axiom by delegating interpretation outside the structure; the paradigm proposed here satisfies it by embedding interpretation within the monitored system itself. Table~\ref{tab:conv-vs-aware} contrasts the two approaches across the dimensions that matter for autonomous assessment.

We define \textbf{Structural Cognition} as the capacity of an instrumented structure to perceive, process, interpret and report its own mechanical state through a permanent, bidirectional coupling between its physical instantiation and a continuously updated digital representation. A \textbf{Self-Aware Steel Structure} is the realisation of this capacity in a steel system: a structure that does not merely transmit raw data but maintains an internal, evolving model of what constitutes its own healthy behaviour, evaluates incoming responses against that model, recognises deviations attributable to stiffness change, connection loosening or progressive damage, and autonomously communicates the resulting assessment. Cognition here is operational rather than metaphorical: it denotes a closed loop in which the structure's measured dynamics are continuously confronted with model-based and data-based expectations, and the residual between expectation and observation is itself the carrier of self-knowledge. This framing is consistent with the data-and-model fusion view of the digital twin, in which predictions are a principled combination of physics models and streaming measurement data \citep{brunton2019data,chiachio2022structural}, and with the physical basis that natural frequencies, mode shapes and damping depend only on mass and stiffness \citep{chopra2017dynamics,cook2002fea}, so that a change in stiffness can produce detectable changes in modal parameters under adequate excitation, sensing and observability conditions.

\subsection{Levels of structural self-awareness}\label{subsec:levels}

To make the paradigm operational and measurable, we propose a six-level model of structural self-awareness, ordered by the depth of internal interpretation a structure performs on its own response. The levels are cumulative: each presupposes the capabilities of those below it. Fig.~\ref{fig:levels} depicts the progression from a passive member to an autonomous agent, and Table~\ref{tab:levels} gives the formal definition of each level together with its mapping onto the Rytter damage-identification hierarchy.

\begin{itemize}
\item \textbf{L0, Passive.} A purely load-resisting member with no instrumentation. Condition is established only by periodic human inspection; the two-state comparison on which damage assessment depends \citep{farrar2013shm} exists, if at all, in the inspector's judgement.
\item \textbf{L1, Sensing.} A distributed sensor network acquires and transmits raw signals (strain, acceleration, temperature, displacement, inclination). This is the level at which most deployed SHM systems terminate: data are stored and transmitted, but the structure performs no interpretation. It corresponds to \textit{data acquisition}, the second part of the statistical pattern recognition paradigm \citep{farrar2013shm}.
\item \textbf{L2, Perception (feature extraction).} The system reduces high-dimensional raw streams to low-dimensional, damage-sensitive features: natural frequencies, mode shapes, damping ratios, dominant modes obtained by SVD/POD, or autoencoder latent variables \citep{farrar2013shm,brunton2019data}. This satisfies Axiom IVa by performing the feature extraction that sensors alone cannot, and it embeds \textit{data normalisation} so that benign operational and environmental variability, notably temperature-induced stiffness change, is separated from response changes caused by damage \citep{farrar2013shm,figueiredo2010machine}.
\item \textbf{L3, Interpretation (diagnosis).} The structure decides whether its current behaviour is consistent with its healthy baseline, and, where it is not, localises and classifies the anomaly. Unsupervised novelty/outlier detection, trained on healthy data only, flags abnormal conditions; supervised classifiers (Random Forest, SVM, ANN) assign damage type and level \citep{farrar2013shm,naser2023ml,sohn2003statistical}. Diagnosis is the act of forming an internal verdict about state, minimising false-positive and false-negative indications \citep{farrar2013shm}.
\item \textbf{L4, Prognosis.} The structure estimates the implications of its diagnosed state for future performance and remaining useful life, fusing the diagnosed condition with model-based predictions of how the identified change propagates. This is \textit{damage prognosis} in the SHM sense \citep{farrar2013shm} and depends on a predictive digital twin able to extrapolate beyond the observed condition \citep{zou2026deep}.
\item \textbf{L5, Autonomous reporting and decision.} The structure communicates its self-assessment without prompting and, where authorised, triggers actions: alarms, maintenance requests, load restrictions, or re-calibration of its own baseline. This is the level at which the structure becomes an agent supporting condition-based and predictive maintenance \citep{farrar2013shm,naser2023ml}, closing the loop between physical event and digital decision.
\end{itemize}

\subsection{Mapping onto the Rytter damage-identification hierarchy}\label{subsec:rytter}

The self-awareness levels are deliberately aligned with the established Rytter hierarchy of damage identification, which orders the questions a diagnostic system can answer by increasing knowledge of the damage state: existence, location, type, severity (extent) and prognosis \citep{farrar2013shm}. \textbf{L2} establishes the feature basis on which any of these questions can be posed. \textbf{L3} answers \textit{existence} and \textit{location}, which Axiom III holds can be addressed in an unsupervised mode, and then advances to \textit{type} and \textit{severity}, which Axiom III holds generally require supervised learning \citep{farrar2013shm,wangcha2022unsupervised}. \textbf{L4} corresponds directly to \textit{prognosis}, the estimation of remaining safe or useful life. \textbf{L5} adds no new diagnostic question but operationalises the answers by communicating and acting upon them autonomously. The hierarchy therefore supplies the diagnostic semantics of structural cognition, while the levels supply its architectural and behavioural semantics: the first describes \textit{what} is known about the damage, the second \textit{how autonomously the structure itself comes to know and act upon it}.

\subsection{The perceive--process--interpret--report loop and the physical--digital link}\label{subsec:loop}

The transition from L1 to L5 is not a sequence of discrete products but a continuous loop executed in place. The \textbf{perceive} stage corresponds to L1 sensing across the heterogeneous, distributed network, with data fusion across same-type arrays (for example, strain gauges) and across heterogeneous channels (acceleration, temperature, inclination) to improve fidelity \citep{farrar2013shm}. The \textbf{process} stage (L2) compresses and normalises these streams into stable features. The \textbf{interpret} stage (L3 and L4) confronts the features with the structure's internal expectation of itself, a finite-element-derived modal and harmonic baseline \citep{chopra2017dynamics,cook2002fea} continuously corrected by data through the digital twin \citep{brunton2019data}, and converts residuals into diagnosis and prognosis. The \textbf{report} stage (L5) returns the verdict to operators and, where permitted, to the structure's own control of baselines and alarms.

What distinguishes this loop from conventional SHM is the permanence and bidirectionality of the physical--digital link. In conventional practice the link is intermittent and one-directional: data flow from structure to archive, and human interpretation flows back only at inspection intervals. In a Self-Aware Steel Structure the link is unbroken, since the digital twin is continuously fed by streaming data and continuously returns updated expectations and verdicts \citep{brunton2019data,li2025edge}, so that interpretation becomes a property the structure carries with it rather than a service performed upon it. The fundamental premise that damage requires a comparison between two system states \citep{farrar2013shm} is thereby internalised: the healthy baseline ceases to be an archived record and becomes a living, model-and-data representation against which the structure perpetually measures itself. This is an enrichment of SHM, not a replacement: every level rests on the statistical pattern recognition paradigm and the SHM axioms \citep{farrar2013shm}. Structural Cognition is the claim that these well-established processes, when embedded, closed and made autonomous, transform a passive member into a structure that understands, and reports, its own condition.
